\begin{table}[!t]
\caption{Table 1: Effect of Russian Invasion on Media Sentiment (RDD, Day Running Variable)} 
\fontsize{12.0pt}{14.4pt}\selectfont
\begin{tabular*}{\linewidth}{@{\extracolsep{\fill}}lccccccc}
\toprule
  & (1) & (2) & (3) & (4) & (5) & (6) & (7) \\ 
\midrule\addlinespace[2.5pt]
Invasion & -0.165*** & -0.154*** & -0.131*** & -0.117*** & -0.267*** & -0.205*** & -0.27*** \\ 
{} & {(0.021)} & {(0.017)} & {(0.013)} & {(0.012)} & {(0.034)} & {(0.033)} & {(0.04)} \\ 
Topic FEs & - & - & - & X & X & X & X \\ 
State FEs & - & X & X & X & - & - & - \\ 
Source FEs & - & - & X & X & - & - & - \\ 
Source Z-Score & - & - & - & - & X & X & X \\ 
Clusters & 61 & 61 & 61 & 61 & 61 & 61 & 61 \\ 
Num.Obs. & 174094 & 174094 & 174094 & 174094 & 174056 & 130488 & 93364 \\ 
\bottomrule
\end{tabular*}
\begin{minipage}{\linewidth}
+ p < 0.1, * p < 0.05, ** p < 0.01, *** p < 0.001\\
All standard errors are clustered at the source level. Models 1-4 use the
optimal bandwidth. Models 5-7 use a Z-Score adjustment instead of fixed effects.
Model 6 omits Military Topics. Model 7 omits NATO allies.\\
\end{minipage}
\end{table}

